% =============================================================================
% TERMO DE CONSENTIMENTO COMUNITÁRIO
% Modelo de Ata de Consentimento Coletivo
% =============================================================================

\documentclass[12pt,a4paper,brazil]{article}

\usepackage[utf8]{inputenc}
\usepackage[T1]{fontenc}
\usepackage[brazil]{babel}
\usepackage{geometry}
\usepackage{setspace}
\usepackage{indentfirst}
\usepackage{graphicx}
\usepackage{fancyhdr}
\usepackage{longtable}
\usepackage{array}
\usepackage{booktabs}

\geometry{a4paper, left=3cm, right=2cm, top=3cm, bottom=2cm}
\onehalfspacing

\pagestyle{fancy}
\fancyhf{}
\rhead{\small Termo de Consentimento Comunitário}
\lhead{\small PPGPI/UFS}
\cfoot{\thepage}

\begin{document}

\begin{center}
{\large\bfseries UNIVERSIDADE FEDERAL DE SERGIPE}\\[0.3cm]
{\normalsize Programa de Pós-Graduação em Ciência da Propriedade Intelectual (PPGPI)}\\[1cm]
{\Large\bfseries TERMO DE CONSENTIMENTO COMUNITÁRIO}\\[0.3cm]
{\normalsize (Ata de Consentimento Coletivo Prévio, Livre e Informado)}\\[0.5cm]
\noindent\rule{\textwidth}{0.4pt}
\end{center}

\vspace{0.5cm}

\noindent\textbf{Título da Pesquisa.} Modelo Computacional para Salvaguarda de Saberes Agrícolas Tradicionais como Ativos Intangíveis em Comunidades Quilombolas do Semiárido Nordeste~II.

\noindent\textbf{Pesquisadora Responsável.} Catuxe Varjão de Santana Oliveira (PPGPI/UFS).

\noindent\textbf{Comunidade Quilombola.} \rule{8cm}{0.4pt}

\noindent\textbf{Município.} Jeremoabo --- BA.

\noindent\textbf{Data da Reunião.} \rule{2cm}{0.4pt} / \rule{2cm}{0.4pt} / \rule{3cm}{0.4pt}

\noindent\textbf{Local da Reunião.} \rule{8cm}{0.4pt}

\vspace{0.5cm}

\section*{Registro da Reunião}

Aos \rule{2cm}{0.4pt} dias do mês de \rule{3cm}{0.4pt} de \rule{3cm}{0.4pt}, na sede da \rule{6cm}{0.4pt} (associação/espaço comunitário), reuniram-se membros da comunidade quilombola \rule{6cm}{0.4pt}, certificada pela Fundação Cultural Palmares, juntamente com a pesquisadora Catuxe Varjão de Santana Oliveira, para apresentação e deliberação sobre a realização da pesquisa acima identificada no território da comunidade.

\section*{Apresentação da Pesquisa}

A pesquisadora apresentou de forma clara e em linguagem acessível os seguintes aspectos da pesquisa.

\paragraph{Objetivos.} A pesquisa visa desenvolver e validar uma ferramenta computacional que auxilie as comunidades quilombolas do Território de Identidade Semiárido Nordeste~II na proteção dos seus próprios saberes agrícolas tradicionais. Para isso, será adaptado um questionário internacional para a linguagem e a realidade das comunidades, com a participação direta de agricultores e mestres de saberes, e será criado um índice que ajuda a identificar quais saberes podem precisar de atenção mais urgente. A ferramenta ficará à disposição das comunidades e poderá ser utilizada em processos junto a IPHAN, INPI, CGen e outros órgãos, fortalecendo a capacidade das comunidades de proteger e valorizar seus conhecimentos por conta própria.

\paragraph{Atividades previstas na comunidade.} As atividades incluem a aplicação de questionários sobre práticas agrícolas tradicionais (em formato de entrevista assistida), entrevistas individuais gravadas em áudio, e oficinas devolutivas para apresentação e validação dos resultados pela comunidade.

\paragraph{Participação.} A participação de cada membro da comunidade será voluntária e individual, mediante assinatura de Termo de Consentimento Livre e Esclarecido (TCLE). Nenhuma pessoa será obrigada a participar, e qualquer participante poderá desistir a qualquer momento sem justificativa e sem prejuízo.

\paragraph{Proteção dos saberes.} Os conhecimentos tradicionais compartilhados serão tratados em conformidade com a Lei nº~13.123/2015 (Marco Legal da Biodiversidade) e com o Protocolo de Nagoia. A pesquisa não visa apropriar-se dos saberes da comunidade, mas sim criar ferramentas de documentação e proteção que fiquem à disposição da própria comunidade. Nenhum conhecimento específico será divulgado sem autorização.

\paragraph{Devolutiva dos resultados.} Ao final da pesquisa, a comunidade receberá relatório simplificado com os resultados referentes ao seu território, incluindo perfil de vulnerabilidade e recomendações de salvaguarda. Os mestres de saberes que participarem do comitê de especialistas serão reconhecidos como coautores do instrumento adaptado.

\paragraph{Riscos e mitigação.} Os riscos são mínimos. A equipe de pesquisa respeitará os calendários agrícolas e as festividades comunitárias para agendamento das atividades. Os dados serão anonimizados e apresentados de forma agregada. A comunidade poderá indicar facilitadores locais para auxiliar na mediação cultural durante a coleta.

\section*{Deliberação}

Após a apresentação, foi aberto espaço para perguntas e esclarecimentos dos membros da comunidade. Todas as dúvidas foram respondidas pela pesquisadora.

\vspace{0.5cm}

\noindent Submetida à deliberação dos presentes, a comunidade decidiu, de forma livre e informada, pela

\vspace{0.3cm}

\noindent
$\square$ \textbf{Autorização} da realização da pesquisa no território da comunidade, nos termos apresentados.

\noindent
$\square$ \textbf{Autorização condicionada}, com as seguintes ressalvas e condições adicionais.

\noindent
\rule{\textwidth}{0.4pt}\\
\rule{\textwidth}{0.4pt}\\
\rule{\textwidth}{0.4pt}

\noindent
$\square$ \textbf{Não autorização} da realização da pesquisa, pelos seguintes motivos.

\noindent
\rule{\textwidth}{0.4pt}\\
\rule{\textwidth}{0.4pt}

\section*{Lista de Presença}

\begin{longtable}{|p{0.5cm}|p{5.5cm}|p{3cm}|p{4cm}|}
\hline
\textbf{Nº} & \textbf{Nome completo} & \textbf{Função na comunidade} & \textbf{Assinatura} \\
\hline
1 & & & \\[0.8cm]
\hline
2 & & & \\[0.8cm]
\hline
3 & & & \\[0.8cm]
\hline
4 & & & \\[0.8cm]
\hline
5 & & & \\[0.8cm]
\hline
6 & & & \\[0.8cm]
\hline
7 & & & \\[0.8cm]
\hline
8 & & & \\[0.8cm]
\hline
9 & & & \\[0.8cm]
\hline
10 & & & \\[0.8cm]
\hline
11 & & & \\[0.8cm]
\hline
12 & & & \\[0.8cm]
\hline
13 & & & \\[0.8cm]
\hline
14 & & & \\[0.8cm]
\hline
15 & & & \\[0.8cm]
\hline
\end{longtable}

\vspace{0.5cm}

\noindent Nada mais havendo a tratar, foi encerrada a reunião, da qual se lavrou a presente ata, que vai assinada pela pesquisadora responsável e pelas lideranças comunitárias presentes.

\vspace{2cm}

\noindent
\begin{minipage}[t]{0.48\textwidth}
\centering
\rule{7cm}{0.4pt}\\
Catuxe Varjão de Santana Oliveira\\
Pesquisadora Responsável --- PPGPI/UFS
\end{minipage}
\hfill
\begin{minipage}[t]{0.48\textwidth}
\centering
\rule{7cm}{0.4pt}\\
Presidente da Associação Comunitária\\
Nome: \rule{5cm}{0.4pt}
\end{minipage}

\vspace{2cm}

\noindent
\begin{minipage}[t]{0.48\textwidth}
\centering
\rule{7cm}{0.4pt}\\
Liderança Comunitária\\
Nome: \rule{5cm}{0.4pt}
\end{minipage}
\hfill
\begin{minipage}[t]{0.48\textwidth}
\centering
\rule{7cm}{0.4pt}\\
Liderança Comunitária\\
Nome: \rule{5cm}{0.4pt}
\end{minipage}

\vspace{1cm}

\noindent
\footnotesize{Este documento é emitido em duas vias de igual teor, ficando uma com a associação comunitária e outra com a pesquisadora responsável. Uma cópia será anexada ao processo junto ao CEP/UFS.}

\end{document}
